\section{Introduction}
\label{sec:intro}
Metaheuristics that learn while they search are a prominent line of work in evolutionary computation and operations research. A typical \emph{learnheuristic}~\cite{CalvetEtal2017} embeds a reinforcement-learning agent, most often tabular Q-Learning~\cite{WatkinsDayan1992,SuttonBarto2018}, as an adaptive selector of operators, strategies or parameters inside a search framework such as Iterated Local Search (ILS)~\cite{LourencaoEtal2019} or Iterated Greedy (IG)~\cite{RuizStutzle2007}. The usual evidence is a comparison with the search without the agent, or with published baselines~\cite{KarimiMamaghanEtal2023,MissaouiOSullivan2026}; when the learned version wins, the result is read as evidence that the \emph{learned policy} is what helps.

That reading is fragile for three reasons. First, the \emph{search structure} that hosts the agent (its operators, its acceptance rule, its local search) may produce most of the gain under any selection rule. Second, the \emph{budget}: an agent changes the CPU time per iteration, so equal iteration counts are not equal budgets, and a policy that merely chooses cheaper actions buys extra iterations without having decided anything better. Third, the \emph{tuning}: a parameter validated offline on separate instances is a competitor that most studies omit, although a practitioner would use it before reaching for a learner. Adaptive operator selection (AOS) is well studied~\cite{FialhoEtal2010,PisingerRopke2007}, but in an audit of $23$ recent open-access papers that embed a learned selector in a metaheuristic, $10$ include a random-selection control, $6$ compare at a common time limit, $3$ include a static configuration tuned offline, and $22$ report a benefit of the learned component ($5$ of them only partly). The audit is small and has a single coder; its coding sheet is in the repository.

\paragraph{Research questions} \textbf{RQ1:} what explains the advantage of a learning-guided IG over ILS, and how much of a local search's gain is only budget? \textbf{RQ2:} do online learners beat random choice over the same actions, at equal iterations and at equal CPU time; does a static configuration win anyway; and could learning have shown up if it existed? \textbf{RQ3:} where, if anywhere, does an online selector pay, can that be foretold, and what exploits the headroom instead?

\paragraph{Contributions} (i) A \emph{budget-fair attribution protocol} (Section~\ref{sec:protocol}): a random control at a matched budget, a split of a learner's gain at equal CPU time into a decision effect and a cost effect, a static competitor validated offline, a cross-fitted oracle headroom with a formal treatment of its bias, power checks, positive controls, an order-sensitivity check, a map of the offline knowledge given to the static competitor, and predictions recorded with code hashes before every battery; the protocol is released as a small library. (ii) Evidence from $\nbatteries$ pre-specified batteries on six hosts (identical parallel machines, unrelated machines, a permutation flowshop, a genetic algorithm and, outside scheduling, an iterated local search for the travelling salesman problem and an adaptive large neighbourhood search for vehicle routing built on a third-party library) that the advantage of the learning-guided IG comes from its search structure and that the standard Q-Learning wins at equal CPU time by choosing cheaper actions, not by a learned policy. (iii) A \emph{headroom analysis}: the usual one-seed oracle headroom is mostly run-to-run noise; with the noise removed it is a necessary condition for a Q-Learning that does not lock in to pay (rank correlation $0.81$ across $28$ classes), never shown sufficient, and a policy learned \emph{offline} exploits it. (iv) A \emph{map of offline knowledge} and the finding that an online learner is best used as a \emph{corrector} of offline knowledge: warm-started with an offline prior it adds value where the prior was transferred and costs where it was exact, on a scheduling benchmark and, against a deliberately wrong prior, on the TSP. (v) A candid account of two reversals: the standard Q-Learning seemed to recover half to two thirds of the validated action's gain on published instances, but it repeats the action listed first; and our first order-invariant control, a permutation of the action list once per run, gains in none of $85$ cells but is itself a lock-in learner: with ties broken at random the learner gains beyond the margin in $32$ of the $85$ cells, in step with the headroom, and recovers $12$--$18\%$ of the validated gain on the standard benchmark. (vi) The price of offline knowledge: tuning a validated action costs $11$ times its test run, one tuning instance per class gives $81\%$ of its gain, and with the tuning charged its advantage over random choice depends on how many instances each class will see.

\paragraph{Outline} Section~\ref{sec:related} places the work among parameter control, algorithm selection and configuration, and benchmarking methodology. Section~\ref{sec:method} describes the hosts, learners and protocol. Section~\ref{sec:results} answers the research questions and reports the order diagnostic. Section~\ref{sec:external} tests how far the answers carry: published instances, equal CPU time, a longer horizon, offline policies, the value of offline knowledge and a second domain. Section~\ref{sec:discussion} discusses when online learning pays, gives a checklist and lists the limitations. This is the condensed version of the full manuscript, whose supplementary sections hold every table, battery and proof (the numbers here are checked against them by an audit script).
